\section[Cote 115 : la dualité d'Isbell par le calcul des noyaux]{Cote 115 : la dualité d'Isbell par le calcul des noyaux\protect\verifie{Folder115}}\label{sec:115}

La cote 115, \emph{\q{Correspondances} fonctorielles. Dualité des topos}, quatorze pages que l'inventaire date \q{[à partir de 1982]}, attache, aux pages~3 à~5, à un noyau $X \in (A \times B)^\wedge$, préfaisceau contravariant en ses deux variables sur deux petites catégories, deux transformées ${}^sX : \widehat A^{\mathrm{op}} \to \widehat B$ et ${}^dX : \widehat B^{\mathrm{op}} \to \widehat A$, reliées par $\mathrm{Hom}_{\widehat A}(L, {}^dX(M)) \simeq X(L, M) \simeq \mathrm{Hom}_{\widehat B}(M, {}^sX(L))$. La page~7 prend $B = A^{\mathrm{op}}$ et pour noyau le hom de $A$, que la marge de la page~3 désigne comme l'élément canonique de $(A \times A^{\mathrm{op}})^\wedge$ (\lot{115}{1}{3}). La lecture modernisée note $\mathcal O$ et $\mathrm{Spec}$ les foncteurs que le manuscrit écrit $\varphi_A$ et $\Psi_A$.

\begin{theoreme}[Grothendieck, pages 3 à 7, la dualité d'Isbell comme cas particulier\piste{115-isbell-duality-derived-from-kernels}]
Pour tout noyau $K$ sur $A \times B$, ${}^sK$, vu comme foncteur $\widehat A \to (\widehat B)^{\mathrm{op}}$, est adjoint à gauche de ${}^dK$. Pour $B = A^{\mathrm{op}}$ et $K = \mathrm{Hom}_A$, les deux transformées sont exactement $\mathcal O(\Phi)(a) = \mathrm{Hom}(\Phi, y a)$ et $\mathrm{Spec}(\Psi)(a) = \mathrm{Hom}(\Psi, \mathrm{Hom}_A(a, -))$, et l'adjonction est $\mathcal O \dashv \mathrm{Spec} : \widehat A \rightleftarrows (A^\vee)^{\mathrm{op}}$. On a $\mathcal O(y a) \simeq \mathrm{Hom}_A(a, -)$, $\mathrm{Spec}(\mathrm{Hom}_A(a, -)) \simeq y a$, et $\mathcal O$ transforme les colimites de $\widehat A$ en limites de $A^\vee$.
\end{theoreme}

L'énoncé tient tel que la lecture le donne. En Lean, le noyau est un foncteur $K : A^{\mathrm{op}} \to [C, \mathbf{Ens}]$, le cas de la page étant $C = B^{\mathrm{op}}$ ; l'adjonction générale est \texttt{kernelAdj}, la dualité \texttt{isbell}, obtenue en posant $K$ égal au foncteur de co-Yoneda, et les deux formules de la page~7 valent par définition (\texttt{isbellO\_obj\_obj}, \texttt{isbellSpec\_obj\_obj}). La question de la piste, savoir si une démonstration publiée obtient ainsi la dualité, relève de la littérature et reste ouverte.

Ne sont pas démontrés ici : la roue des onze catégories de foncteurs de la page~5, l'extension aux familles finies $(A_i)$, la caractérisation de $\mathcal O$ et de $\mathrm{Spec}$ par leur restriction le long de $y$ et la dualité des topos des pages suivantes.
